% ls-many-authors.tex -- Chapter 4 addition (writer: authors).
% Inserted after the section labelled sec:consent-canon, before the
% unnumbered "Review of the key ideas". Everything below is unbuilt.

\section{Many authors, one sovereign}\label{sec:ls-many-authors}

Until now one person has authorized the swarm. Section~\ref{sec:roles} read
that requirement off the grant lifecycle: the four steps need ``exactly one
participant who originates authority and can withdraw it.'' Bob's repository
already breaks the premise in a small way. Bob owns the code and its
acceptance test; Alice runs the team. Nothing stops Bob from directing an
agent himself, and nothing stops a third person from joining. This section
asks what survives when $M$ people direct $N$ agents through one daemon,
what the single-operator design must add, and which conflicts between their
directions can be found before any agent acts.

The answer reuses the kernel, the grant, the commitment checker and the claim.
What it adds is an owner table, a precedence rule, and a routing rule for
escalations. Nothing in this section exists in the reference implementation.
Each rule below states its maturity; experiments are protocols, not results.

\subsection{Three direct answers}\label{sec:ls-many-authors-answers}

\paragraph{Does only one person have administrative access?} No, but there is
one charter. The single-writer guarantees concern the daemon, not the people
who address it. Every mutation still passes through one writer
(\pdchapref{swk}{The Single-Writer Kernel}%
\ifpdbook
, Definition~\ref{swk:def:single-writer}%
\fi
), so the exclusion and fencing results and the synchronous-guard result hold
for any number of requesters%
\ifpdbook
\ (Theorems~\ref{swk:thm:exclusion} and~\ref{swk:thm:decidability})%
\fi
. Their requests serialize like anyone else's. What changes is the author
slot. Several people now originate authority, so the quoted premise is
false, and the escalation lane of \S\ref{sec:escalation} no longer knows whom
to interrupt. The smallest repair is not one administrator. It is one
\emph{charter}: a root object that records who delegated what to whom, which
rule settles a clash, and how the charter itself may change.

\paragraph{Can one person control all the agents?} In authority, yes. Every grant descends from the charter,
so whoever holds the charter can withdraw all of them through the override
(Invariant~\ref{prop:override}). In attention, not in general.
Theorem~\ref{thm:escalation-threshold} shows that the feasible debit band can
be empty at a given attention budget. Every principal whose conflicts route to
one person adds to that person's queue. The incident command system faces the
same limit with people. Its doctrine sets an optimal span of control of one
supervisor to five subordinates~\pdcite{dhs2017nims}. That figure is a
guideline for people supervising people; nothing here transfers it to agents.

\paragraph{What happens when directions conflict?} Detection and resolution
separate. Detection is mechanical wherever directions fit the commitment
fragment of \pdchapref{he}{The Harbor Economy}: polynomial, with a witness,
and NP-complete once a direction may offer a choice. For program behavior in
general it is undecidable. Resolution is normative. That chapter already says
so: ``no consistency checker chooses their normative content.'' A recent
benchmark shows why the distinction matters. On tasks whose two
specifications edit the same code incompatibly, one agent given both
specifications solved $62\%$, while two agents given one specification each
solved $16\%$ (Claude Code with Claude Opus~4.8). The merged code typically
satisfied only one specification, and a communication channel between the
agents recovered only a small part of the gap~\pdcite{mediratta2026parallact}.
The two requests were jointly satisfiable. What was missing was a party
entitled to decide how they fit together.

\subsection{Nine rules for many authors}\label{sec:ls-many-authors-rules}

The rules below are the smallest set we found that keeps one sovereign while
admitting several authors. A rule marked \Designed{} is specified completely
here and reuses a check the book already specifies. A rule marked \Vision{}
needs policy content or mechanisms not yet written down.

\begin{pdclaim}{Design invariant}{Many-Author Authority}
\setlength{\parskip}{4pt plus 1pt}
\label{def:ls-many-authors}
A harbor directed by principals $1,\dots,M$ satisfies:
\begin{enumerate}[leftmargin=2.2em,itemsep=2pt,label=R\arabic*.]
\item \textbf{One charter root} (\Vision). One charter is the root of every
delegation and holds the precedence table. An amendment may require $k$ of $M$
signatures; each amendment is one proposal, accepted or rejected.
\item \textbf{A delegation tree} (\Designed). After transitive reduction, every
scope, invariant and agent grant has exactly one delegating parent. The
daemon refuses a grant that would give a node a second parent.
\item \textbf{Attenuation per person} (\Vision). The charter's allocation to
principal $i$ bounds every grant $i$ issues, and each agent card attenuates the
grant of the principal it acts for%
\ifpdbook
\ (Invariant~\ref{anchor:def:offline-attn})%
\fi
.
\item \textbf{One principal per action} (\Designed). Every admitted action
names exactly one grant. The witness of the consent-grant lifecycle
(\S\ref{sec:consent}) becomes $(p, g, x, \text{reason}, \text{artifact})$,
with $p$ the issuing principal. An agent may hold grants from two people; the
guard checks each action against the one grant it names, never their union.
\item \textbf{One owner per scope and per joint invariant} (\Vision). This
includes an \emph{integration owner} for ``the combined head builds and the
preserved tests pass.'' Owner precedence is a default the owner may waive.
\item \textbf{Directions become tagged commitments} (\Designed). Each
direction compiles to commitment sentences tagged with their issuer, stored
beside the verbatim words. Admission runs the tagged checker over the union of
accepted commitments%
\ifpdbook
\ (Theorem~\ref{he:thm:he-tagged-blame})%
\fi
\ and against the declared reachable futures.
\item \textbf{Precedence} (\Vision). Between principals, the higher delegator
prevails, then the owner of the clashing scope. Within one principal, a later
direction supersedes an earlier one. Time alone never decides between two
principals.
\item \textbf{Escalation routing} (\Vision). A conflict that precedence does
not settle goes, with its witness, to the least common authority of the scopes
it touches (Theorem~\ref{thm:ls-unity-of-command}), as one proposal to accept
or reject.
\item \textbf{Exit per principal} (\Vision). Each principal's override is
unconditional over that principal's own grants. Revoking another principal's
grant needs authority above it. A principal who does not own a shared scope
may pause only its safety-class actions, for a declared time-to-live, and the
pause escalates to the scope's owner.
\end{enumerate}
\end{pdclaim}

\noindent Three of these rules answer objections the single-operator design
never met.

Rule R1 is not ``one administrator'' under another name. The incident command
system's \emph{unified command} lets several agencies set ``a common set of
incident objectives, strategies, and a single'' incident action plan, and the
arrangement ``does not affect individual agency authority, responsibility, or
accountability''~\pdcite{dhs2017nims}. The charter plays that part: one plan,
several retained authorities. Ostrom's polycentric governance, ``many centers
of decision making that are formally independent of each
other''~\pdcite{ostrom2010polycentric}, describes the relation between harbors,
which \pdchapref{fh}{The Federated Harbor} develops. Inside one harbor the
centers are not independent. They share a writer.

Rule R4 is the capability answer to an old failure. Hardy's compiler wrote its
billing file on a user's behalf because, in his words, ``the compiler runs with
authority stemming from two sources''~\pdcite{hardy1988confused}. An agent that
works for Alice and Bob is in the compiler's position. Naming one grant per
action is the fix capabilities supplied: the action states which authority it
exercises, so it cannot spend Bob's authority on Alice's request. Grants
themselves are already capabilities with caveats~\pdcite{macaroons2014}.

Rule R7 deliberately refuses ``last direction wins'' between people. The norm
literature's three heuristics are lex specialis (the more specific norm), lex
posterior (the later norm) and lex superior (the higher
authority)~\pdcite{olson2024defeasible}. Olson and colleagues derive lex
posterior inside a single normative theory and leave merging across
individuals to future work. Between two people, lateness measures only who
spoke second.

Rule R8 has a price. Routing every escalation in a scope to its one owner is
the sole-owner arm of the queueing boundary (Theorem~\ref{thm:specialization});
letting any of several owners answer is the pooled arm. GitHub's code-owner
review takes the pooled arm: ``an approval from any of the owners is
sufficient''~\pdcite{githubcodeowners}. Theorem~\ref{thm:specialization} prices
that choice. It is not free in either direction.

Rule R9 bounds the override. An unconditional global stop held by each of $M$
people gives the swarm $M$ independent ways to halt. Exit stays inalienable
over one's own grants. Over the shared work it becomes a bounded pause with an
owner who must answer it.

\subsection{Which conflicts can be found before an agent acts}\label{sec:ls-many-authors-conflicts}

Directions can be compared at three moments. At \emph{instruction time} the
daemon holds only commitments: what each person obliged, forbade, claimed and
scheduled. At \emph{claim time} it also holds what each agent plans to touch.
At \emph{admission time} it holds the actual combined state that a commit would
produce. Each moment sees more and comes later.
Table~\ref{tab:ls-conflict-classes} sorts conflicts by the earliest moment that
can see them. ``Detectable'' means a stated procedure decides the question; it
does not mean the procedure has been run on this system.

\begin{table}[tp]
\centering
\caption{Conflict classes between directions from several principals, the
earliest moment each can be found, and what happens when it is. Rows 2--5 use
the commitment checker of \pdchapref{he}{The Harbor Economy}; rows 6--8 concern
program behavior, where detection weakens and then fails.}
\label{tab:ls-conflict-classes}
\small
\renewcommand{\arraystretch}{1.15}
\begin{tabularx}{\linewidth}{@{}>{\raggedright\arraybackslash}p{0.22\linewidth} >{\raggedright\arraybackslash}X >{\raggedright\arraybackslash}X >{\raggedright\arraybackslash}X@{}}
\toprule
\textbf{Class (example)} & \textbf{Found before execution?} & \textbf{Mechanism} & \textbf{When found} \\
\midrule
1.\ Claim overlap (two agents claim one file) & Yes, at claim time; complete for the overlap relation & Overlap scan; fenced lease & Release, expiry, bonded preemption, then a human ruling routed by R8 \\
2.\ Opposed directions, including derived ones (migrate vs.\ freeze) & Yes, at instruction time; polynomial, with a witness & Tagged commitment checker & Precedence (R7); otherwise escalate (R8) \\
3.\ Choices that cannot all be met (``use A or B'') & Decidable but NP-complete, even when each person alone is consistent & Bounded search for a checkable selection & A timeout refuses admission; escalate with the partial witness \\
4.\ Infeasible schedule (deadline vs.\ review window) & Yes; a negative cycle in the deadline graph & Shortest-path feasibility test & Owner of the deadline \\
5.\ A later direction fires a dormant rule & Only against a declared family of reachable futures & Recheck at every admission & Later issuer informed; owner resolves \\
6.\ Clean textual merge, broken build (signature change vs.\ caller) & Partly at claim time, by symbol claims and a dependency graph; certainly only by building the combined state & Speculative merge and build~\pdcite{brun2011proactive}; build-graph conflict analysis~\pdcite{ananthanarayanan2019keeping}; combined-head guard & Integration owner issues a narrow joint grant \\
7.\ Merged code meets one specification only & No: undecidable in general~\pdcite{horwitz1989integrating} & Safe approximation; tests as a partial oracle & Both specification owners; dissent recorded \\
8.\ Two people editing one text live & Convergence is guaranteed; intent is not~\pdcite{shapiro2011crdt,kleppmann2019interleaving} & CRDT or operational transform, for text & Result treated as a proposal for admission \\
9.\ Disagreement about values, with no fact to find & Detectable as a clash; not decidable by computation & None; the charter decides & One binary proposal to the scope owner \\
\bottomrule
\end{tabularx}
\end{table}

Three rows deserve a sentence. Row~1's completeness is real but narrow. The
advisory-claim theorem of \pdchapref{bonded}{The Bonded Commons} reports every
overlap%
\ifpdbook
\ (Theorem~\ref{bonded:thm:advisory})%
\fi
, and its overlap relation is about files. A signature change in one file and a
call in another overlap nowhere, so row~6 lies outside that guarantee
entirely. Horwitz, Prins and Reps state row~7's limit precisely: ``it is
undecidable whether a program modification actually leads to'' a change in
behavior, so their integration algorithm computes ``a safe
approximation''~\pdcite{horwitz1989integrating}. Row~8 is the collaborative
editor's version of the same gap. Strong eventual consistency guarantees that
replicas which have delivered the same updates have equivalent
state~\pdcite{shapiro2011crdt}, and Kleppmann and colleagues show that this
property ``is necessary but it is not sufficient'' for a usable
editor~\pdcite{kleppmann2019interleaving}. Convergence is a property of the
text, not of what either person meant by it.

\subsection{Unity of command and the escalation path}\label{sec:ls-many-authors-escalation}

Rule R8 needs a target. The incident command system's answer is \emph{unity of
command}: ``each individual only reports to one person''~\pdcite{dhs2017nims}.
The theorem below states exactly when that answer is needed for escalation to
be well defined.

Model delegation as a finite directed acyclic graph $G$. Its nodes are
principals and \emph{objects} (scopes, invariants and agent grants). An edge
$u\to v$ means principal $u$ delegated authority to $v$. Objects have no
outgoing edges, the charter root $r$ has no incoming edge, and every node is
reachable from $r$. A principal $p$ \emph{holds authority over} object $s$ when
$G$ has a path from $p$ to $s$. A conflict witness touches a nonempty set $S$
of objects. Its \emph{common authorities} are
\[
C(S)=\{\,p : p \text{ holds authority over every } s\in S\,\},
\]
which contains $r$. Escalation for $S$ is \emph{well defined} when $C(S)$ is
totally ordered by reachability. Its lowest element is then the first hop, and
each further hop is the next element up to $r$.

\begin{pdclaim}{Theorem}{Unity of Command Is Exactly Well-Defined Escalation}
\setlength{\parskip}{4pt plus 1pt}
\label{thm:ls-unity-of-command}
Escalation is well defined for every nonempty set of objects if and only if,
in the transitive reduction of $G$, every object and every principal with a
path to an object has at most one parent.
\end{pdclaim}
\begin{proof}
\setlength{\parskip}{4pt plus 1pt}
Write $A(v)$ for the set of principals with a path to $v$. The transitive
reduction of a finite acyclic graph has the same paths between the same pairs
of nodes, so $A(v)$ may be read off the reduction.

If every such node has at most one parent, follow parents upward from an
object $s$. Each step is forced, and every principal with a path to $s$ lies on
this one route to $r$, so $A(s)$ is a chain. For a set $S$, $C(S)$ is the
intersection of the chains $A(s)$, $s\in S$. It is a subset of one chain, hence
a chain, and it contains $r$.

Conversely, let $v$ be an object $s$, or a principal with a path to an
object $s$, and let $v$ have two parents $p\ne q$ in the transitive reduction.
Both are principals, since only principals delegate. Neither reaches the
other: if $p$ reached $q$, the path from $p$ through $q$ to $v$ would imply the
edge from $p$ to $v$, and the reduction would have removed it; symmetrically
for $q$. Both reach $s$ through $v$, so both lie in $C(\{s\})$, which is
therefore not a chain.
\end{proof}

The theorem is elementary; its use is to name the design rule. A tie-break
that picks one parent per node restores a unique path, and it does so by
choosing a spanning tree: the rule is unity of command again, chosen by an
ordering instead of by assignment. GitHub's code-owner file shows the tree-making move in
practice: when several patterns match a file, ``the last matching pattern
takes the most precedence''~\pdcite{githubcodeowners}. Order turns overlapping
patterns into one owning line per file.

\subsection{Aggregation: where Arrow applies}\label{sec:ls-many-authors-aggregation}

Chapter~0 applied Arrow's theorem to agents voting on proposals. In a
many-author harbor the aggregation problem belongs to the people.

The admission rule above aggregates by conjunction: a plan is admissible when
it satisfies every accepted commitment of every principal. That rule is not a
dictatorship, but it is partial. When the commitments cannot all hold, it
refuses to choose, and something else must. Suppose that something must rank
at least three joint plans completely, for every profile of the principals'
rankings, while respecting unanimous preferences and independence of
irrelevant alternatives. Arrow's theorem, as stated in Chapter~0, says it is
then a dictatorship. The Gibbard--Satterthwaite theorem, stated in the same
place, adds that a deterministic single-winner rule over at least three
outcomes is dictatorial or open to manipulation at some profile. Fickinger and colleagues carry the point to
assistants serving several humans. When the humans' own demonstrations count
for little against what the assistant then does, the only straightforward
deterministic mechanisms are those restricted to two outcomes and
dictatorships~\pdcite{fickinger2020multi}.

The design consequence is the shape of rules R1, R5 and R8. Escalations are put
as one proposal to accept or reject, which keeps every decision to two
outcomes. Each scope has one decider, the local dictator these theorems make
unavoidable for complete ranked choice. Nothing here requires a vote, and the
design holds none.

Owner decisiveness has its own limit. Giving each owner decisive authority
over its scope, together with the Pareto requirement, is the conflict
\pdchapref{bonded}{The Bonded Commons} draws from Sen's theorem for agents'
locks%
\ifpdbook
\ (Theorem~\ref{bonded:prop:sen})%
\fi
~\pdcite{sen1970impossibility}. With people in place of agents the warning
reads the same way: an owner who cannot be asked to yield can block an outcome
everyone prefers. That is why R5 makes precedence waivable, and why the
direct-release step of that chapter's dispute ladder comes first. Money does
not dissolve the problem either. When several principals pay one agent,
Bernheim and Whinston find the aggregate incentive scheme cost-minimizing, but
noncooperative principals induce an efficient action only if collusion among
them would implement the first-best action at first-best
cost~\pdcite{bernheim1986common}. ``Account, don't auction''
(\S\ref{sec:account}) was a single-operator rule. With several operators it
becomes an open question.

\begin{pdexample}{Two Authors, Three Agents, One Afternoon}
\setlength{\parskip}{4pt plus 1pt}
Bob holds the charter. His grant to Alice covers writes in \texttt{payments/},
reads in \texttt{tests/checkout/}, and preservation of the original checkout
test, for four hours. Bob keeps \texttt{db:checkout}, writes in
\texttt{tests/}, and the integration invariant. Heron and Wren act under
Alice's grants; Pike acts under Bob's. Times are minutes from now. This is a
constructed run, not a measurement.

\textbf{Instruction time: an opposed direction.} Alice tells Wren to store
amounts as integer cents. It compiles to the fact \texttt{cents\_requested},
the Horn rule \texttt{cents\_requested}~$\to$~\texttt{schema\_change}, and
Alice's rule $\texttt{schema\_change}\to
\mathsf O(\texttt{migrate},\{\texttt{db:checkout}\},[0,120])$. Bob's standing
release freeze is $\mathsf F(\texttt{migrate},\{\texttt{db:checkout}\},
[0,4320])$, three days being $3\times24\times60=4320$ minutes. One Horn step
fires the obligation, and the intervals overlap on
$[\max(0,0),\min(120,4320)]=[0,120]$: a conflict. The deletion filter makes two
checker calls. Without Alice the set holds no obligation, so no conflict, and
Alice stays. Without Bob it holds no prohibition, and Bob stays. The blame set
is $\{\text{Alice},\text{Bob}\}$. The scope \texttt{db:checkout} is Bob's, so
precedence (R7) keeps the freeze. The daemon refuses Wren's plan before Wren
starts, and Alice receives the witness. Nobody had to decide anything.

\textbf{Instruction time: an infeasible schedule.} Alice asks Bob for one
exception. Bob will allow the migration if he reviews it. Let $s$ be the
migration's start, $m$ its end and $e$ the end of Bob's review. Wren
estimates thirty minutes, so $m-s\ge30$, that is $s-m\le-30$. Bob's rule needs
sixty minutes of review, so $e-m\ge60$, that is $m-e\le-60$. Alice's demo
needs it done within eighty minutes: $e-s\le80$. Summing the three
inequalities, the left side cancels to $0$ and the right side is
\[
-30-60+80=-10,
\]
so $0\le-10$ is required. The cycle $s\to e\to m\to s$ has weight $-10$, and
the schedule is infeasible. Each constraint is needed, so the blame set is
again both people. The witness touches Alice's demo and Bob's review rule; the
only common authority is Bob. He receives one proposal: shorten the review to
fifty minutes? The cycle weight becomes $-30-50+80=0$, which is feasible, with
no slack. Bob accepts.

\textbf{Claim time: a warning.} Heron plans to change \texttt{charge(amount)}
to \texttt{charge(amount\_cents)} in \texttt{payments/}. Pike plans a refund
test in \texttt{tests/checkout/} that calls \texttt{charge}. The two file
claims do not overlap, so the advisory-claim theorem correctly reports
nothing. A symbol-level read claim by Pike against Heron's modify claim on the
same function would raise a warning (\Vision).

\textbf{Admission time: a refusal.} Pike's test lands first, under Bob's grant.
Heron's patch reaches the guard, which builds the combined head. The refund
test calls the old signature, the build fails, and admission is refused with
both commits as the witness. Neither owner can repair it alone: the call is in
Bob's \texttt{tests/}, the signature in Alice's \texttt{payments/}. As
integration owner, Bob issues Alice a narrow grant: \texttt{tests/checkout/},
symbol \texttt{charge}, thirty minutes. Alice attenuates it onto Heron's card,
and Heron's fix names that grant (R4).

\textbf{Tally.} Two conflicts were refused before any agent worked, one was
flagged at claim time, and one was refused at admission. No broken commit
landed. Bob decided twice and Alice never did. The load fell on the root,
which is the queueing concern of \S\ref{sec:sole-roles}.

\emph{Now you try:} find the longest review that fits Alice's demo, and redo it
for a seventy-minute demo, in Exercise~\ref{ex:ls-many-authors-cycle}.
\end{pdexample}

\subsection{Disagreements about facts are bets}\label{sec:ls-many-authors-bets}

Not every conflict is about values. Alice may believe the migration is safe
while Bob believes it will break staging. Critch and Russell show that an
agent serving principals with different beliefs cannot, in general, be
Pareto-optimal by maximizing a fixed weighted sum of their utilities. A
Pareto-optimal policy must shift weight toward the principal whose prior better
predicts what the agent observes~\pdcite{critch2017servant}. Their example is a
cake. Alice is $90\%$ sure it is red, Bob $90\%$ sure it is green, and each
values half a cake at $20$ and a whole cake at $30$. The policy ``all to Alice
if red, all to Bob if green'' gives each an expected $0.9\times30+0.1\times0
=27$ by their own beliefs, against $20$ for an even split, and no fixed
weighting implements it. The result assumes the principals know each other's
priors.

The harbor should not let an agent learn such shifting weights silently.
Authority that drifts toward whoever predicted better is a change of authority
nobody consented to. The proposal (\Vision) is to make the bet explicit, as a
float plan of \pdchapref{he}{The Harbor Economy}%
\ifpdbook
\ (\S\ref{he:sec:float-plan})%
\fi
: each person posts a bond against the observable outcome, the terms appear in
the digest, and settlement follows the outcome, not the agent's inference.
Exercise~\ref{ex:ls-many-authors-bets} works the cake by hand.

\begin{pdboundary}
\setlength{\parskip}{4pt plus 1pt}
Compiling a natural-language direction into commitments is itself an
interpretive step, and nothing here verifies it. Storing the verbatim words
beside the compiled sentences preserves what was said; it does not prove the
compilation faithful.

Theorem~\ref{thm:ls-unity-of-command} is elementary graph structure, and the
aggregation argument applies theorems proved elsewhere. Neither shows that a
tree of owners is the right organization for a given team, only that a
well-defined escalation path requires one.

The benchmark figures come from a diagnostic set built to require
coordination; they say how often agents fail when such conflicts arise, not how
often the conflicts arise. Who ought to hold the charter, and whether an
owner's decision is legitimate, are questions for the people involved.
\end{pdboundary}

\pdexercisepointer{\ref{ex:ls-many-authors-tree}}{\ref{ex:ls-many-authors-pause}}{\pageref{ex:ls-many-authors-tree}}
